\documentclass[10pt,a4paper]{amsart}
\usepackage[margin=19mm]{geometry}
\usepackage{amsmath,amssymb,mathtools}
\usepackage[scr=rsfs,cal=euler]{mathalfa}
\usepackage{xfrac}
\usepackage[unicode,hidelinks,bookmarks=false]{hyperref}
\newcommand{\ie}{i.e.\ }
\newcommand{\cf}{cf.\ }
\newcommand{\resp}{respectively\ }
\newcommand{\Subsection}{\subsection}
\providecommand{\nobreakdash}{\nobreak\hbox{-}}
\setlength{\emergencystretch}{2em}
\multlinegap0pt
\usepackage{bm}
\usepackage{bbm}
\usepackage{dsfont}
\usepackage{xspace}
\usepackage{cancel}
\usepackage{mathtools}
\usepackage{amsthm}
\makeatletter
\@ifundefined{thm}{\newtheorem{thm}{Thm}}{}
\@ifundefined{lem}{\newtheorem{lem}[thm]{Lemma}}{}
\makeatother
\newcommand{\HSI}{\textsf{HSI}}
\newcommand{\up}[1]{\textup{#1}}
\title[Multiplicatively idempotent HSI algebras satisfy all equatio]{Multiplicatively idempotent HSI algebras satisfy all equations of $\mathbb{N}$}
\author{Tumadhir Alsulami, Marcel Jackson, Michael Kinyon}
\date{Source: arXiv:2609.16922v1 (CC BY 4.0)}
\begin{document}
\maketitle

\noindent\emph{Source and licence.} Multiplicatively idempotent HSI algebras satisfy all equations of $\mathbb{N}$. Version arXiv:2609.16922v1 (CC BY 4.0), distributed under the Creative Commons Attribution 4.0 International licence, as recorded in the arXiv licence metadata for this exact version and shown on the abstract page https://arxiv.org/abs/2609.16922v1. Full rights notice: licenses/arxiv-2609.16922-CC-BY-4.0.txt.\par\smallskip

\noindent\textit{Editorial scope.} Counted unit: the Lem whose source label is \texttt{lem:canincreasemod2} in the article (source0.tex lines 630--632, source label \texttt{lem:canincreasemod2}), delivered together with the article's own complete original proof of it (source0.tex lines 633--652, one \texttt{proof} environment of 20 lines). No mathematics is written, repaired or re-derived: statement and proof are the article's own text line for line. Required context retained uncounted: lem:x2x (lines 565--573); lem:equivmod2 (lines 616--618). Every definition, display and label of those lines is retained unchanged. Omitted from this package and not redistributed: 1--564, 574--615, 619--629, 653--1937. Nothing inside a retained excerpt is omitted or altered: every retained body is delivered exactly as the article's lines are written, so no omission receipt of any kind accompanies this package. Citations: the retained excerpts carry no citation command at all, so no bibliography entry of the article is transcribed and no reference is redistributed. Reference adaptation: none was needed and none was made; every reference inside the delivered unit resolves inside it, so no unresolved reference or citation is produced. Printed slips kept literal: no printed word, symbol, hypothesis or number of the counted statement or of its proof was altered. Layout and class: the article's own class is \texttt{amsart} with packages including amssymb, latexsym, amsmath, amsthm, enumerate, mathabx, eucal, framed, color, graphicx, mathrsfs, fontenc, xy, tikz; the wrapper is the lane's pinned \texttt{amsart} editorial wrapper at 10pt on A4, which restates the theorem-like environments the retained text needs and adds the article's own macro definitions that the retained text needs (\texttt{\textbackslash HSI}, \texttt{\textbackslash up}), with extra wrapper packages (bm, bbm, dsfont, xspace, cancel, mathtools). The delivered numbering is the unit's own. Assets: no retained span contains a figure environment, a graphics-inclusion command, a PDF-page inclusion, a table or a TikZ picture; no image file is copied into this package, the package declares no asset dependency and no image file is redistributed with it. Licence: version arXiv:2609.16922v1 of this article is distributed under the Creative Commons Attribution 4.0 International licence, as recorded in the arXiv licence metadata for this exact version and shown on the abstract page https://arxiv.org/abs/2609.16922v1. A full rights notice is delivered at licenses/arxiv-2609.16922-CC-BY-4.0.txt. Characters: every retained excerpt is pure ASCII, so the delivered engine, XeLaTeX with its default Latin Modern text font, reports no missing character.
\begin{lem}\label{lem:x2x}
The following laws are consequences of $\HSI\cup\{x^2\approx x\}$\up:
\begin{align*}
x+y\approx {}&x+2xy +y,\\
1+x\approx {}& 1+3x,\\
x+xy\approx{}& x+3xy,\\
2\approx{}& 4.
\end{align*}
\end{lem}

\begin{lem}\label{lem:equivmod2}
Every monomial appearing in $s$ with coefficient $c\geq 1$ appears in $s'$ with coefficient $c'\geq c$ and $c\equiv_2 c'$.
\end{lem}

\begin{lem}\label{lem:canincreasemod2}
If a monomial $m$ appearing in $s'$ is not strictly covered by some monomials in $s$, then either it has the same coefficient in $s$ as in $s'$ or  $\HSI\cup\{x^2\approx x\}\vdash s\approx s+2m$.
\end{lem}

\begin{proof}
Let $X=\{x_1,\dots,x_u\}$ denote the variables in $m$ and $Y=\{y_1,\dots,y_v\}$ the variables in $s$ that are not in $m$.
First we note that $m$ must be at least covered by monomials in $s$.  
To show this, let $Y$ denote the variables not appearing in $m$, and consider the law $s|_{Y=0}\approx s'|_{Y=0}$.  
Now $m$ appears in $s'|_{Y=0}$ and the same variables appear in  $s|_{Y=0}$. 
It follows that all of the variables in $m$ appear in monomials of $s$ whose content is a subset of that in $m$.
Equivalently, $m$ is covered by the monomials in $s|_{Y=0}$ and hence in $s$.

As $m$ is covered but not strictly covered, it follows that there is at least one monomial appearing in $s$ with the same variables as $m$.  
If a monomial with the same variables as $m$ (including possibly $m$ itself) appears in~$s$ with coefficient at least $2$, then the result follows using $2\approx 4$ and Lemma \ref{lem:equivmod2}.  
Similarly if there is more than one monomial $m_1,m_2$ in $s$ with the same variables as $m$, then $\HSI\cup\{x^2\approx x\}\vdash m_1\approx m_2\approx m$ and then $2\approx 4$ can again be used to obtain $s\approx s+2m$.   
Now assume the remaining case: there is a unique monomial $m_s$ in $s$ with the same variables as $m$, and that the coefficient of $m_s$ in $s$ is $1$.   
If there is a monomial $m'$ in $s$ whose variables are a proper subset of those in $m$ and $m_s$, then the law $x+xy\approx x+3xy$ from Lemma~\ref{lem:x2x} gives $m'+m_s\approx m'+3m_s$ and we are in a case already considered.  Note that it is allowed for $m'$ to be a constant (with set of variables $\varnothing$).  
So we can now assume that $m_s$ is the only monomial in $s$ whose variables are a subset of $X$ and the assumptions on $m_s$ in $s$ then imply that $s|_{Y=0}$ is simply $m_s$ by itself.  Our goal is to show that the coefficient of $m$ in $s'$ is also $1$.    
The factorisation of $s$ as $\prod_{1\leq i\leq k}(p_i-q_i)$ continues to hold after assigning variables in $Y$ to $0$, so each $p_i-q_i$ has $(p_i-q_i)|_{Y=0}$ equal to a factor of $m_s$: either a constant, or a monomial whose variables are within $X$.  
We denote the factor $(p_i-q_i)|_{Y=0}$ as $m_i$, and observe that the reduction of $p_i-q_i$ to $(p_i-q_i)|_{Y=0}$ is simply by replacing monomials with a variable in $Y$ by $0$, there are no proper subtractions that arise between nonzero coefficient monomials $m'-m'$, as these did not originally occur in $p_i-q_i$.
This in turn implies that the coefficient of $m_i$ cannot be negative in $(p_i-q_i)|_{Y=0}$, for otherwise in $p_i-q_i$, all monomials in $p_i$ involved at least one variable in $Y$: taking these values very close to $0$ with values for $X$ variables fixed then makes $p_i-q_i$ negative (yet it is a positive polynomial).
It follows that $m_i$ is the only monomial in $p_i$ whose variables are within $X$ (including the possibility that $m_i$ is a constant), and that all monomials in $q_i$ involve a variable from $Y$. 
Thus the unique monomial in either of $s=\prod_{1\leq i\leq k}(p_i-q_i)$ (which we denoted $m_s$) or $s'=\prod_{1\leq i\leq k}(p_i+q_i)$ (which must be $m$) with variables amongst $X$ is the monomial $m_1m_2\dots m_k$, and the coefficient is fixed identically in both cases.
\end{proof}

\end{document}
